% GENERATED FILE. Edit canonical lessons, then run npm run generate:books.
% canonical-source-hash: fnv1a64:8891a3e10dd7c28a
% canonical-lessons: ZH-C192-cesuo, ZH-C192-dianti, ZH-C192-tingchechang, ZH-C192-paidui, ZH-C192-kaifang

\chapter{Signs in a Building}
\label{ch:zh-signs-in-a-building}
\hlchaptermodality{192}

\begin{chapteropening}
\textbf{By the end of this chapter:} \emph{I can ask for the toilets, the lift and the car park, queue, and say when a place is open.}

Complete the last lesson of chapter 192: I can ask for the toilets, the lift and the car park, queue, and say when a place is open.
\end{chapteropening}

\section[cèsuǒ]{cèsuǒ — toilet; lavatory}

\label{lesson:ZH-C192-cesuo}



\begin{quote}
Before the new one: say the Mandarin for to go abroad, then the Mandarin for to travel for work.
\end{quote}

\subsection*{You'll want to know: cèsuǒ}
\emph{cèsuǒ} — \textquotedblleft{}toilet; lavatory\textquotedblright{}.

\begin{quote}
\emph{qǐngwèn cèsuǒ zài nǎr} — Excuse me, where's the toilet?
\end{quote}

Its characters come later; for now it is a word for the ear.

\begin{cousinweb}
\emph{cèsuǒ} is built from \emph{cè} 'a privy' and \emph{suǒ} 'a place', both inherited from Middle Chinese, \emph{tsrhiH} and \emph{srjoX} (Guangyun, in Baxter's transcription). The modern \emph{cè} is not the regular outcome of that form, and how it came is not settled.
\end{cousinweb}

\begin{grammarlens}[title={What you've built}]
One more word for reading notices and messages.
\end{grammarlens}

\subsection*{Your turn}
\begin{itemize}
\raggedright
  \item \emph{Say it:} \emph{cèsuǒ}
  \item \emph{Say it:} \emph{cèsuǒ}, once more
  \item \emph{Say it:} \emph{qǐngwèn cèsuǒ zài nǎr} — Excuse me, where's the toilet?
  \item \emph{Recall:} say the Mandarin for to go abroad, then the Mandarin for to travel for work, then say \emph{cèsuǒ} again
\end{itemize}

\begin{tcolorbox}[breakable,colback=teal!4,colframe=teal!35!black,title={Before you move on}]
What does cèsuǒ mean? (\textquotedblleft{}Toilet; lavatory\textquotedblright{}.) Say it once more.
\end{tcolorbox}
\section[diàntī]{diàntī — lift; elevator}

\label{lesson:ZH-C192-dianti}



\begin{quote}
Before the new one: say the Mandarin for to travel for work, then the Mandarin for lavatory.
\end{quote}

\subsection*{You'll want to know: diàntī}
\emph{diàntī} — \textquotedblleft{}lift; elevator\textquotedblright{}.

\begin{quote}
\emph{diàntī huài le} — The lift is out of order.
\end{quote}

Its characters come later; for now it is a word for the ear.

\begin{cousinweb}
\emph{diàntī} is built from \zh{电} 'electricity' and \emph{tī} 'a ladder', the \emph{tī} of \emph{lóutī}, both inherited from Middle Chinese, \emph{denH} and \emph{thej} (Guangyun, in Baxter's transcription), which glosses the second as wooden steps.
\end{cousinweb}

\begin{grammarlens}[title={What you've built}]
One more word for reading notices and messages.
\end{grammarlens}

\subsection*{Your turn}
\begin{itemize}
\raggedright
  \item \emph{Say it:} \emph{diàntī}
  \item \emph{Say it:} \emph{diàntī}, once more
  \item \emph{Say it:} \emph{diàntī huài le} — The lift is out of order.
  \item \emph{Recall:} say the Mandarin for to travel for work, then the Mandarin for lavatory, then say \emph{diàntī} again
\end{itemize}

\begin{tcolorbox}[breakable,colback=teal!4,colframe=teal!35!black,title={Before you move on}]
What does diàntī mean? (\textquotedblleft{}Lift; elevator\textquotedblright{}.) Say it once more.
\end{tcolorbox}
\section[tíngchēchǎng]{tíngchēchǎng — car park}

\label{lesson:ZH-C192-tingchechang}



\begin{quote}
Before the new one: say the Mandarin for lavatory, then the Mandarin for lift (in a building).
\end{quote}

\subsection*{You'll want to know: tíngchēchǎng}
\emph{tíngchēchǎng} — \textquotedblleft{}car park\textquotedblright{}.

\begin{quote}
\emph{shāngchǎng xiàbian yǒu yí ge tíngchēchǎng} — There's a car park under the shopping centre.
\end{quote}

Its characters come later; for now it is a word for the ear.

\begin{cousinweb}
\emph{tíngchēchǎng} is built from \emph{tíng} 'to stop', \zh{车} 'vehicle' and \emph{chǎng} 'an open ground', all inherited from Middle Chinese, \emph{deng}, \emph{tsyhae} and \emph{drjang} (Guangyun, in Baxter's transcription).
\end{cousinweb}

\begin{grammarlens}[title={What you've built}]
One more word for reading notices and messages.
\end{grammarlens}

\subsection*{Your turn}
\begin{itemize}
\raggedright
  \item \emph{Say it:} \emph{tíngchēchǎng}
  \item \emph{Say it:} \emph{tíngchēchǎng}, once more
  \item \emph{Say it:} \emph{shāngchǎng xiàbian yǒu yí ge tíngchēchǎng} — There's a car park under the shopping centre.
  \item \emph{Recall:} say the Mandarin for lavatory, then the Mandarin for lift (in a building), then say \emph{tíngchēchǎng} again
\end{itemize}

\begin{tcolorbox}[breakable,colback=teal!4,colframe=teal!35!black,title={Before you move on}]
What does tíngchēchǎng mean? (\textquotedblleft{}Car park\textquotedblright{}.) Say it once more.
\end{tcolorbox}
\section[páiduì]{páiduì — to queue; to line up}

\label{lesson:ZH-C192-paidui}



\begin{quote}
Before the new one: say the Mandarin for lift (in a building), then the Mandarin for car park.
\end{quote}

\subsection*{You'll want to know: páiduì}
\emph{páiduì} — \textquotedblleft{}to queue; to line up\textquotedblright{}.

\begin{quote}
\emph{dàjiā dōu zài páiduì} — Everyone is queuing.
\end{quote}

Its characters come later; for now it is a word for the ear.

\begin{cousinweb}
\emph{páiduì} is built from \emph{pái} 'to set in a row' and \emph{duì} 'a file of people'. The second is inherited from Middle Chinese \emph{dwojH} (Guangyun, in Baxter's transcription), glossed 'a troop'; the dictionary glosses the first 'to push aside', so how it came to mean lining up is not settled.
\end{cousinweb}

\begin{grammarlens}[title={What you've built}]
One more word for reading notices and messages.
\end{grammarlens}

\subsection*{Your turn}
\begin{itemize}
\raggedright
  \item \emph{Say it:} \emph{páiduì}
  \item \emph{Say it:} \emph{páiduì}, once more
  \item \emph{Say it:} \emph{dàjiā dōu zài páiduì} — Everyone is queuing.
  \item \emph{Recall:} say the Mandarin for lift (in a building), then the Mandarin for car park, then say \emph{páiduì} again
\end{itemize}

\begin{tcolorbox}[breakable,colback=teal!4,colframe=teal!35!black,title={Before you move on}]
What does páiduì mean? (\textquotedblleft{}To queue; to line up\textquotedblright{}.) Say it once more.
\end{tcolorbox}
\section[kāifàng]{kāifàng — to be open to the public}

\label{lesson:ZH-C192-kaifang}



\begin{quote}
Before the new one: say the Mandarin for car park, then the Mandarin for to queue.
\end{quote}

\subsection*{You'll want to know: kāifàng}
\emph{kāifàng} — \textquotedblleft{}to be open to the public\textquotedblright{}.

\begin{quote}
\emph{gōngyuán wǎnshang bù kāifàng} — The park isn't open at night.
\end{quote}

Its characters come later; for now it is a word for the ear.

\begin{cousinweb}
\emph{kāifàng} is built from \zh{开} 'to open' and \emph{fàng} 'to let go', both inherited from Middle Chinese, \emph{khoj} and \emph{pjangH} (Guangyun, in Baxter's transcription).
\end{cousinweb}

\begin{grammarlens}[title={What you've built}]
One more word for reading notices and messages.
\end{grammarlens}

\subsection*{Your turn}
\begin{itemize}
\raggedright
  \item \emph{Say it:} \emph{kāifàng}
  \item \emph{Say it:} \emph{kāifàng}, once more
  \item \emph{Say it:} \emph{gōngyuán wǎnshang bù kāifàng} — The park isn't open at night.
  \item \emph{Recall:} say the Mandarin for car park, then the Mandarin for to queue, then say \emph{kāifàng} again
\end{itemize}

\begin{tcolorbox}[breakable,colback=teal!4,colframe=teal!35!black,title={Before you move on}]
What does kāifàng mean? (\textquotedblleft{}To be open to the public\textquotedblright{}.) Say it once more.
\end{tcolorbox}
